\par\addvspace{12pt}\Needspace{9\baselineskip}

\section{仿射变换在圆锥曲线中的应用}

% ITEM: quadratic-example-035
\solutionitem{quadratic-example-035}{例4.35}
设\(E:2x^2+2xy+y^2=1\)，\(A=(1,-1)\)。点\(B\)在\(E\)上运动且\(B\ne A\)，\(M\)是\(AB\)的中点。选择变换将\(E\)化为圆，求\(M\)的轨迹，并将所得结论还原为原坐标中的方程。
\solutionanswer
取\(u=x,v=x+y\)，则\(E\)成为单位圆，\(A\)映为\((1,0)\)。圆中弦中点\((s,t)\)满足\((s-1/2)^2+t^2=1/4\)，去掉\((1,0)\)。还原\(s=X,t=X+Y\)，轨迹为\((X-1/2)^2+(X+Y)^2=1/4\)，去掉\((1,-1)\)。逆向取\(B=2M-A\)可恢复每个允许中点。


% ITEM: quadratic-exercise-23-01
\solutionitem{quadratic-exercise-23-01}{习题4.23（1）}
将\((x-1)^2/9+(y+2)^2/4=1\)化为单位圆，求点\((4,-2)\)及直线\(2x-3y=5\)的像。求原图中三角形面积与像三角形面积的关系。
\par
\solutionanswer
取圆化变换\(u=(x-1)/3,v=(y+2)/2\)，椭圆变为\(u^2+v^2=1\)。点\((4,-2)\)的像为\((1,0)\)。直线像由\(x=3u+1,y=2v-2\)代入得
\[2(3u+1)-3(2v-2)=5,\quad 6u-6v+3=0.\]
此变换面积倍率为\(1/(3\cdot2)=1/6\)，故任意原图三角形的像三角形面积等于原面积的\(1/6\)，即原面积为像面积的6倍。


% ITEM: quadratic-exercise-23-02
\solutionitem{quadratic-exercise-23-02}{习题4.23（2）}
将\(x^2+4xy+5y^2=1\)化为单位圆，写出逆变换，并求该椭圆的面积。
\par
\solutionanswer
配方为\(x^2+4xy+5y^2=(x+2y)^2+y^2\)。取\(u=x+2y,v=y\)，则椭圆化为单位圆\(u^2+v^2=1\)，逆变换为\(x=u-2v,y=v\)。逆变换线性部分行列式为1，面积倍率为1，故原椭圆面积为\(\pi\)。


% ITEM: quadratic-exercise-23-03
\solutionitem{quadratic-exercise-23-03}{习题4.23（3）}
对椭圆\(2(x-1)^2+2(x-1)(y+1)+(y+1)^2=1\)，选择变换化为单位圆，求中心及面积，并写出直线\(x+y=0\)的像。
\par
\solutionanswer
取\(u=x-1,v=x+y\)，则二次式成为\(u^2+v^2\)，所以化为单位圆。逆变换\(x=u+1,y=v-u-1\)给出中心\((1,-1)\)，面积倍率为\(1\)，原面积为\(\pi\)。直线\(x+y=0\)映为\(v=0\)。


% ITEM: quadratic-example-036
\solutionitem{quadratic-example-036}{例4.36}
在椭圆\(x^2/9+y^2/4=1\)中，求与\(y=x\)平行的弦的中点所在直线，并确定中点的实际取值范围。若中点给定为\(M(1,-4/9)\)，求相应弦所在直线。
\solutionanswer
平行弦中点满足\(4x+9y=0\)。令\(M=(t,-4t/9)\)，中点在椭圆内部当且仅当\(13t^2/81<1\)，故\(|t|<9/\sqrt{13}\)。给定中点时，弦线为过\((1,-4/9)\)且斜率\(1\)的直线\(y=x-13/9\)。


% ITEM: quadratic-example-037
\solutionitem{quadratic-example-037}{例4.37}
设\(A,B\)是椭圆\(x^2/9+y^2/4=1\)上一条固定弦的两端点，\(A=(3,0)\)、\(B=(0,2)\)，\(P\)在椭圆上运动。求\(\triangle ABP\)的重心\(G\)的轨迹。若要求三角形非退化，轨迹应怎样取？
\solutionanswer
由\(G=(A+B+P)/3\)，轨迹为
\[
(x-1)^2+\frac{(y-2/3)^2}{4/9}=1.
\]
非退化条件排除\(P=A,B\)，故轨迹去掉\((2,2/3),(1,4/3)\)。


% ITEM: quadratic-example-038
\solutionitem{quadratic-example-038}{例4.38}
设\(P=(1,1/2)\)，过\(P\)的两条直线分别交椭圆\(x^2/4+y^2=1\)于\(A,C\)和\(B,D\)四个不同的点，且\(AB\parallel CD\)。证明直线\(AB\)的方向固定，求其斜率；再证明过\(P\)且平行于\(AB\)的弦被\(P\)平分。
\solutionanswer
圆化为\((u,v)=(x/2,y)\)，\(P'=(1/2,1/2)\)。圆内接梯形的对角线交点在其对称轴上，该轴经过圆心并垂直两底，故\(A'B'\)方向为\((1,-1)\)。还原后\(AB\)斜率为\(-1/2\)。圆中经过\(P'\)且垂直\(OP'\)的弦被\(P'\)平分，中点性质还原后仍成立，得第二结论。


% ITEM: quadratic-exercise-24-01
\solutionitem{quadratic-exercise-24-01}{习题4.24（1）}
求椭圆\(x^2/16+y^2/9=1\)中与向量\((2,1)\)平行的弦的中点轨迹。
\par
\solutionanswer
先作圆化\((u,v)=(x/4,y/3)\)，弦方向\((2,1)\)变为\((1/2,1/3)\)。圆中半径到弦中点垂直于弦，所以
\[
 \frac{x}{4}\frac12+\frac{y}{3}\frac13=0,
 \qquad 9x+8y=0.
\]
中点须在椭圆内部；将直径方程代入椭圆得\(13x^2/64<1\)。反过来，圆内这条直径上的每一点，都可作为相应方向一条非退化弦的中点；还原后仍成立。因此实际轨迹是
\[
 9x+8y=0,\qquad -\frac8{\sqrt{13}}<x<\frac8{\sqrt{13}}.
\]
另可将端点写成\(M\pm t(2,1)\)，代入原方程相减，直接得到同一直径条件。


% ITEM: quadratic-exercise-24-02
\solutionitem{quadratic-exercise-24-02}{习题4.24（2）}
设\(A=(2,0)\)，\(P\)在椭圆\(x^2/4+y^2=1\)上运动，且\(\overrightarrow{AM}=2\overrightarrow{AP}\)。求点\(M\)的轨迹，再将系数\(2\)改为\(-1\)比较。
\par
\solutionanswer
向量条件给出\(M=2P-A\)，逆向\(P=(M+A)/2\)代入即得\((x+2)^2/16+y^2/4=1\)。系数改为\(-1\)时\(M=2A-P\)，轨迹为\((x-4)^2/4+y^2=1\)。两者中心分别为\((-2,0),(4,0)\)，半轴分别为\(4,2\)与\(2,1\)，对应都可逆，故表示完整轨迹。


% ITEM: quadratic-exercise-24-03
\solutionitem{quadratic-exercise-24-03}{习题4.24（3）}
椭圆\(x^2/9+y^2/4=1\)上两点\(A,B\)满足\(\overrightarrow{PA}=2\overrightarrow{BP}\)，其中\(P=(0,1)\)。求弦\(AB\)所在直线，并用圆化后的几何关系解释结果。
\par
\solutionanswer
令 \(u=x/3,\ v=y/2\)，椭圆变为单位圆，\(P'=(0,1/2)\) 仍按 \(2:1\) 内分弦。设 \(P'B'=s\)，弦中点为 \(H\)，则 \(P'H=s/2,\ HA'=3s/2\)。圆心 \(O\) 到弦的垂线经过 \(H\)，从而
\[
  1-\tfrac14=OA'^2-OP'^2=HA'^2-HP'^2=2s^2 .
\]
\(OP'\) 竖直，它在弦上的投影长为 \(P'H\)。所以弦的倾角 \(\varphi\) 满足 \(|\sin\varphi|=P'H/OP'=s=\sqrt{3/8}\)，即 \(|\tan\varphi|=\sqrt{3/5}\)。还原时斜率乘 \(2/3\)，两条弦线为 \(y=1\pm2x/\sqrt{15}\)。两方向关于纵轴对称，均能取得。


% ITEM: quadratic-exercise-24-04
\solutionitem{quadratic-exercise-24-04}{习题4.24（4）}
（据2013年大纲卷理科第8题）椭圆 \(x^2/4+y^2/3=1\) 的左右顶点为 \(A_1,A_2\)。动点 \(P\) 在椭圆上，若 \(k_{PA_2}\in[-2,-1]\)，求 \(k_{PA_1}\) 的取值范围。若换成任一直径，两斜率之积还相同吗？用圆周角解释。
\par
\solutionanswer
把椭圆化为单位圆：\(u=x/2,\ v=y/\sqrt3\)。直径所对的圆周角为直角，故两条像弦的斜率乘积为 \(-1\)。还原时各斜率乘 \(\sqrt3/2\)，所以 \(k_{PA_1}k_{PA_2}=-3/4\)。由 \(-2\le k_{PA_2}\le-1\) 得 \(3/8\le k_{PA_1}\le3/4\)。圆周角的论证适用于任一直径，故只要两斜率存在，乘积就不变。


% ITEM: quadratic-exercise-24-05
\solutionitem{quadratic-exercise-24-05}{习题4.24（5）}
设\(A,B\)是椭圆\(x^2/9+y^2/4=1\)的一对直径端点，\(P\)在椭圆上且异于\(A,B\)。证明\(\triangle PAB\)的重心轨迹仍为椭圆，求其方程并说明应去掉的点。
\par
\solutionanswer
设直径端点为\(A,B\)，则\(B=-A\)。三角形重心\(G=(A+B+P)/3=P/3\)，所以当\(P\)遍历椭圆时，\(G\)遍历以原点为中心缩小到原来\(1/3\)的椭圆：
\[x_G^2+\frac{9y_G^2}{4}=1.\]
题设排除\(P=A,B\)，相应应从轨迹中去掉\(G=A/3=-B/3\)和\(G=B/3=-A/3\)这两个点。反过来，轨迹上其余任一点\(G\)对应\(P=3G\)为原椭圆上且不同于\(A,B\)的点，故确实由非退化三角形重心取得。


% ITEM: quadratic-example-039
\solutionitem{quadratic-example-039}{例4.39}
（2014年安徽卷理科第14题）设\(F_1,F_2\)分别是椭圆\(E:x^2+y^2/b^2=1\)（\(0<b<1\)）的左、右焦点。过\(F_1\)的直线交椭圆于\(A,B\)，且\(|AF_1|=3|BF_1|\)、\(AF_2\perp x\)轴。求椭圆方程。
\solutionanswer
设\(c=\sqrt{1-b^2}\)，作圆化变换\(u=x,v=y/b\)。两焦点的像为\(F_1'=(-c,0),F_2'=(c,0)\)，这里只把它们作为圆内指定点使用。
由\(AF_2\)竖直及圆方程，\(A'=(c,\pm b)\)，故\(|A'F_1'|^2=4c^2+b^2\)。同一直线上分比保持，有\(|A'F_1'|=3|B'F_1'|\)。圆内相交弦的圆幂关系给出
\[
 |A'F_1'|\,|B'F_1'|=1-c^2=b^2.
\]
因此\(4c^2+b^2=3b^2\)，再结合\(c^2+b^2=1\)，得\(c^2=1/3,b^2=2/3\)。椭圆为\(x^2+3y^2/2=1\)。这条路线同时使用了圆化后的分比、轴向关系与圆幂。

再用焦弦关系：记半通径为\(\ell\)。由\(|AF_1|=3|BF_1|\)及两段焦半径的倒数和为\(2/\ell\)，得\(|AF_1|=2\ell\)。又因\(A\)是右通径的端点，\(|AF_2|=\ell\)。椭圆定义给出\(3\ell=2a\)；而\(a=1,\ \ell=b^2/a\)，故\(b^2=2/3\)。


% ITEM: quadratic-exercise-25-01
\solutionitem{quadratic-exercise-25-01}{习题4.25（1）}
过\(P=(1,0)\)的直线交椭圆\(x^2/9+y^2/4=1\)于\(A,B\)。分别在直线水平、竖直时求\(|PA|\,|PB|\)，再说明这一乘积为什么不是常数。
\par
\solutionanswer
水平弦端点为\((3,0),(-3,0)\)，距离乘积为\(2\cdot4=8\)。竖直弦端点为\((1,\pm4\sqrt2/3)\)，乘积为\(32/9\)。两值不同，故不是常数；圆化保留直线和分比，但不同方向的长度倍率不同。


% ITEM: quadratic-exercise-25-02
\solutionitem{quadratic-exercise-25-02}{习题4.25（2）}
在上题中，\(A=(3,0)\)时，求另一交点\(B\)。再设\(A=(0,2)\)，利用圆化及圆幂求\(|PA|/|PB|\)。
\par
\solutionanswer
第一种情形\(B=(-3,0)\)。第二种将点写为\(P+t(A-P)=(1-t,2t)\)，圆化后的圆幂给出两个参数之积\((q(P)-1)/q(A-P)=-4/5\)。一个参数为\(1\)，另一个为\(-4/5\)，故\(B=(9/5,-8/5)\)，并且\(|PA|/|PB|=5/4\)。参数比例在仿射变换下保持。


% ITEM: quadratic-exercise-25-03
\solutionitem{quadratic-exercise-25-03}{习题4.25（3）}
设\(P\)在椭圆内部，证明过\(P\)的弦\(AB\)满足
\begin{equation*}
    \min(a^2,b^2)\bigl(1-q(P)\bigr)\le |PA|\,|PB|
    \le\max(a^2,b^2)\bigl(1-q(P)\bigr) ,
\end{equation*}
并指出等号对应的方向。
\par
\solutionanswer
取弦的单位方向\(d\)。由根与系数的关系，\(|PA|\,|PB|=(1-q(P))/q(d)\)。单位向量满足\(1/\max(a^2,b^2)\le q(d)\le1/\min(a^2,b^2)\)，相除即得两边界。下界在短轴方向取到，上界在长轴方向取到；圆的情形各方向均取等。


% ITEM: quadratic-exercise-25-04
\solutionitem{quadratic-exercise-25-04}{习题4.25（4）}
过\(M=(1,0)\)的两条不同直线分别交椭圆\(x^2/9+y^2/4=1\)于\(A,B\)及\(C,D\)，斜率分别为\(k_1,k_2\)。若\(|MA|\,|MB|=|MC|\,|MD|\)，证明\(k_1+k_2=0\)。
\par
\solutionanswer
斜率为\(k\)时，方向幂给出距离乘积\(32(1+k^2)/(4+9k^2)\)。令两式相等并交叉相乘，得到\(5(k_1^2-k_2^2)=0\)。两条过同一点的直线不同，所以\(k_1\ne k_2\)，只能\(k_1+k_2=0\)。


% ITEM: quadratic-exercise-25-05
\solutionitem{quadratic-exercise-25-05}{习题4.25（5）}
设\(P=(1,0)\)，过\(P\)的直线交椭圆\(x^2/9+y^2/4=1\)于\(A,B\)。过中心\(O\)作\(OC\parallel AB\)，其中\(C\)在椭圆上。证明\(|PA|\,|PB|/|OC|^2\)为定值，并求此值。
\par
\solutionanswer
将 \((x,y)\) 变为 \((x/3,y/2)\)，椭圆化为单位圆。因 \(PA,PB,OC\) 平行，三段的伸缩倍率相同，在所求比值中恰好抵消。像圆中 \(O'C'=1,\ O'P'=1/3\)，由圆幂即得
\[
 \frac{|PA|\,|PB|}{|OC|^2}=\frac{|P'A'|\,|P'B'|}{|O'C'|^2}=1-(1/3)^2=\frac89.
\]


% ITEM: quadratic-example-040
\solutionitem{quadratic-example-040}{例4.40}
（2013年山东卷文科第22题）椭圆\(C\)的中心在原点，焦点在\(x\)轴上，短轴长为\(2\)，离心率为\(\sqrt2/2\)。


（1）求椭圆方程。
\par



（2）\(A,B\)为椭圆上满足\(S_{\triangle AOB}=\sqrt6/4\)的任意两点，\(E\)为\(AB\)的中点，射线\(OE\)交椭圆于\(P\)。设\(\overrightarrow{OP}=t\overrightarrow{OE}\)，求\(t\)。
\par
\solutionanswer
（1）由\(b=1,e^2=1/2\)，椭圆为\(x^2/2+y^2=1\)。（2）圆化后\(S_{OA'B'}=\sqrt3/4\)，所以两单位半径夹角\(\theta\)满足\(\sin\theta=\sqrt3/2\)，即\(\theta=\pi/3\)或\(2\pi/3\)。中点到圆心距离分别为\(\sqrt3/2,1/2\)，射线与圆交点的半径为\(1\)，故\(t=2/\sqrt3\)或\(2\)。


% ITEM: quadratic-example-041
\solutionitem{quadratic-example-041}{例4.41}
求椭圆\(x^2/9+y^2/4=1\)内接平行四边形的最大面积，分别给出一个边平行于坐标轴和一个边不平行于坐标轴的取等图形。判断所有取得最大面积的图形是否都是矩形。
\solutionanswer
圆化后内接平行四边形最大面积为单位圆内接正方形的面积\(2\)，还原倍率\(6\)，最大\(12\)。一组取等图形的顶点为\((\pm3/\sqrt2,\pm\sqrt2)\)，边平行坐标轴；另一组为\((3,0),(0,2),(-3,0),(0,-2)\)，边不平行坐标轴。这后一图形不是矩形，故所有最大图形不必为矩形。


% ITEM: quadratic-example-042
\solutionitem{quadratic-example-042}{例4.42}
过定点\(T=(0,-2)\)的直线与椭圆\(x^2/4+y^2=1\)交于不同的两点\(A,B\)。求\(S_{\triangle OAB}\)的最大值及取得最大值时直线\(AB\)的方程，其中\(O\)为原点。
\solutionanswer
令 \(u=x/2,\ v=y\)，化椭圆为单位圆。两条单位半径所成三角形的面积至多 \(1/2\)，等号恰在两半径垂直时成立；此时圆心到弦的距离为 \(1/\sqrt2\)。过 \(T=(0,-2)\) 的竖直弦面积为零，其余弦在原图写成 \(y=kx-2\)，像弦为 \(2ku-v=2\)。令其到圆心的距离 \(2/\sqrt{4k^2+1}=1/\sqrt2\)，得 \(k=\pm\sqrt7/2\)。故两条取等直线为 \(y=\pm(\sqrt7/2)x-2\)，原面积最大为 \(2\cdot(1/2)=1\)。


% ITEM: quadratic-exercise-26-01
\solutionitem{quadratic-exercise-26-01}{习题4.26（1）}
设\(A=(-3,0)\)、\(B=(3,0)\)，点\(P\)在椭圆\(x^2/9+y^2/4=1\)上运动。求\(S_{\triangle ABP}\)的最大值，比较直接用高与化椭为圆两种处理。
\par
\solutionanswer
以\(AB=6\)为底，最大高为\(2\)，故最大面积\(6\)，在\(P=(0,\pm2)\)取到。圆化后底是单位圆直径，三角形最大面积为\(1\)，再乘面积倍率\(6\)也得\(6\)。此题直接用高更短。


% ITEM: quadratic-exercise-26-02
\solutionitem{quadratic-exercise-26-02}{习题4.26（2）}
求椭圆\(x^2/a^2+y^2/b^2=1\)中\(S_{\triangle OAB}\)的最大值，并用\(OA,OB\)的共轭关系说明取等位置。
\par
\solutionanswer
圆化后两半径长均为\(1\)，面积至多\(1/2\)，等号对应半径垂直。还原后最大值为\(ab/2\)，取等条件恰为\(OA,OB\)共轭。


% ITEM: quadratic-exercise-26-03
\solutionitem{quadratic-exercise-26-03}{习题4.26（3）}
设\(A=(a,0)\)、\(B=(0,b)\)，点\(P\)在椭圆\(x^2/a^2+y^2/b^2=1\)上运动。求\(S_{\triangle ABP}\)的最大值及对应点\(P\)。
\par
\solutionanswer
圆化后\(A'=(1,0),B'=(0,1)\)，三角形面积为\(|u+v-1|/2\)。由\(|u+v|\le\sqrt2\)，最大值为\((1+\sqrt2)/2\)，在\((u,v)=(-1/\sqrt2,-1/\sqrt2)\)取到。原面积最大\(ab(1+\sqrt2)/2\)，对应\(P=(-a/\sqrt2,-b/\sqrt2)\)。


% ITEM: quadratic-exercise-26-04
\solutionitem{quadratic-exercise-26-04}{习题4.26（4）}
椭圆\(x^2/t^2+y^2/(4-t^2)=1\)，其中\(0<t<2\)。求其内接平行四边形面积的最大可能值，说明为什么不能只求单位圆内接图形的面积而忽略变换倍率。
\par
\solutionanswer
对每个\(t\)，最大面积为\(2t\sqrt{4-t^2}\)。因\(t^2(4-t^2)\le4\)，总最大值为\(4\)，在\(t=\sqrt2\)取到。单位圆中最大面积虽恒为\(2\)，还原倍率\(t\sqrt{4-t^2}\)却随椭圆改变，不能省去。


% ITEM: quadratic-exercise-26-05
\solutionitem{quadratic-exercise-26-05}{习题4.26（5）}
点\(P\)在椭圆\(x^2/9+y^2/4=1\)上，过\(P\)且不过原点\(O\)的直线交椭圆\(x^2/36+y^2/16=1\)于\(A,B\)，射线\(PO\)交外椭圆于\(Q\)。求\(S_{\triangle ABQ}/S_{\triangle ABO}\)及\(S_{\triangle ABQ}\)的最大值。
\par
\solutionanswer
把外椭圆化为单位圆，内椭圆成为半径\(1/2\)的圆。由\(Q=-2P\)及\(AB\)过\(P\)，两三角形到同底\(AB\)的高之比为\(3\)，故面积比为\(3\)。若弦线到圆心距离为\(h\)，则\(0<h\le1/2\)，\(S_{OAB}=h\sqrt{1-h^2}\)，在此区间最大为\(\sqrt3/4\)。还原面积倍率为\(24\)，故\(S_{ABQ}\)最大\(18\sqrt3\)。例如\(P=(3,0),AB:x=3,Q=(-6,0)\)取到。


% ITEM: quadratic-exercise-26-06
\solutionitem{quadratic-exercise-26-06}{习题4.26（6）}
证明椭圆的任意一对共轭直径将其内部区域分成面积相等的四部分，并求\(x^2/9+y^2/4=1\)中每部分的面积。
\par
\solutionanswer
圆化后两共轭直径成为垂直直径，将圆盘分成四个面积相等的扇形。还原时四部分面积均乘同一倍率，故仍相等，每部分为\(6\pi/4=3\pi/2\)。


% ITEM: quadratic-exercise-26-07
\solutionitem{quadratic-exercise-26-07}{习题4.26（7）}
椭圆\(x^2/9+y^2/4=1\)上三点\(A,B,C\)满足\(\overrightarrow{OA}+\overrightarrow{OB}+\overrightarrow{OC}=0\)。证明\(\triangle ABC\)的面积为定值，并求此值。圆化后，这个向量条件对三角形的形状有什么要求？
\par
\solutionanswer
圆化后的三个单位向量\(u,v,w\)满足\(u+v+w=0\)。由\(|u+v|^2=|w|^2=1\)，得\(u\cdot v=-1/2\)，其余两对同理。三点构成正三角形，面积\(3\sqrt3/4\)。还原乘\(6\)，定面积为\(9\sqrt3/2\)。


% ITEM: quadratic-exercise-26-08
\solutionitem{quadratic-exercise-26-08}{习题4.26（8）}
若不要求椭圆\(x^2/a^2+y^2/b^2=1\)的内接四边形为平行四边形，其面积能否超过\(2ab\)？证明你的判断，并确定取得最大值的条件。
\par
\solutionanswer
不能。单位圆内接四边形两对角线长均至多\(2\)，其面积为两对角线乘积与夹角正弦的一半，故不超过\(2\)。等号要求两对角线均为直径且垂直，即四边形为正方形。还原后最大为\(2ab\)，取等当且仅当两对顶点分别是一对共轭直径的端点。


% ITEM: quadratic-example-043
\solutionitem{quadratic-example-043}{例4.43}
（据2019年全国I卷理科第22题改编）曲线\(C\)为椭圆\(x^2+y^2/4=1\)去掉点\((-1,0)\)所得的点集，直线\(l:2x+\sqrt3y+11=0\)。求\(C\)上的点到\(l\)距离的最小值，并确定取得最小值的点。
\solutionanswer
圆化\(x=u,y=2v\)后，\(2x+\sqrt3y=2u+2\sqrt3v\)的最小值为\(-4\)，对应\((u,v)=(-1/2,-\sqrt3/2)\)。因此距离最小为\((11-4)/\sqrt7=\sqrt7\)，取等点为\((-1/2,-\sqrt3)\)，没有被源题排除。


% ITEM: quadratic-example-044
\solutionitem{quadratic-example-044}{例4.44}
在椭圆\(x^2/9+y^2/4=1\)上，\(A=(3,0)\)，点\(B\ne A\)满足\(OA\perp AB\)是否可能？若将条件改为\(OB\perp AB\)，求所有这样的点\(B\)。分别用原坐标与圆化坐标处理，比较需要转化的条件。
\solutionanswer
先在原坐标中计算。\(OA\perp AB\)要求\(AB\)竖直，直线\(x=3\)只有切点\(A\)，故不可能。\(OB\perp AB\)则给出\(x^2+y^2-3x=0\)，与椭圆联立得\(5x^2-27x+36=0\)。排除\(A\)，其余为\(B=(12/5,\pm6/5)\)。

再圆化为\((u,v)=(x/3,y/2)\)，此时\(A'=(1,0)\)、\(B'=(u,v)\)、\(u^2+v^2=1\)。原数量积应写成\(9\Delta u\,\delta u+4\Delta v\,\delta v\)。因此第一种垂直条件为\(9(u-1)=0\)，仍只能给\(B'=A'\)；第二种为
\[
 9u(u-1)+4v^2=0
 \quad\Longrightarrow\quad (u-1)(5u-4)=0.
\]
排除\(u=1\)，得\(u=4/5,v=\pm3/5\)，还原得到同一组\(B\)。圆方程变简单了，但原图的垂直须用加权数量积表达，不能直接要求\(OB'\perp A'B'\)。


% ITEM: quadratic-exercise-27-01
\solutionitem{quadratic-exercise-27-01}{习题4.27（1）}
求\(x^2/9+y^2/4\le1\)中\(2x-y\)的最大值、最小值及取等点。
\par
\solutionanswer
在椭圆盘中，线性式\(2x-y\)的最大、最小值分别为\(\pm\sqrt{2^2\cdot3^2+(-1)^2\cdot2^2}=\pm2\sqrt{10}\)。取等点为
\[
P_{\pm}=\pm\left(\frac{9}{\sqrt{10}},-\frac{2}{\sqrt{10}}\right).
\]
两点满足\(x^2/9+y^2/4=81/90+4/40=1\)，且\(2x-y=\pm20/\sqrt{10}=\pm2\sqrt{10}\)，故确在边界并达到所列极值。


% ITEM: quadratic-exercise-27-02
\solutionitem{quadratic-exercise-27-02}{习题4.27（2）}
求椭圆\(x^2/9+y^2/4=1\)到直线\(x+y=5\)的最近距离与最远距离。若直线改为\(x+y=1\)，最近距离怎样变化？
\par
\solutionanswer
椭圆上\(x+y\)的值域为\([-\sqrt{13},\sqrt{13}]\)。直线\(x+y=5\)在椭圆外，最近、最远距离分别为
\[
\frac{5-\sqrt{13}}{\sqrt2},\qquad \frac{5+\sqrt{13}}{\sqrt2}.
\]
若直线改为\(x+y=1\)，由于\(1\)落在上述值域内，直线与椭圆相交，最近距离变为\(0\)。


% ITEM: quadratic-exercise-27-03
\solutionitem{quadratic-exercise-27-03}{习题4.27（3）}
（2013年上海卷秋季文科第18题）设\(\Omega_n\)为\(x^2/4+ny^2/(4n+1)\le1\)确定的区域，\(M_n\)为\(x+y\)在\(\Omega_n\)上的最大值。求\(M_n\)及\(\lim_{n\to\infty}M_n\)。
\par
\solutionanswer
令 \(x=2u,\ y=\sqrt{4+1/n}\,v\)，区域化为单位圆盘。于是 \(x+y=(2,\sqrt{4+1/n})\cdot(u,v)\) 的最大值就是前一向量的长度 \(M_n=\sqrt{8+1/n}\)，在同向的单位向量处取得。因此 \(\lim_{n\to\infty}M_n=2\sqrt2\)。


% ITEM: quadratic-exercise-27-04
\solutionitem{quadratic-exercise-27-04}{习题4.27（4）}
椭圆\(x^2/9+y^2/4=1\)的两条共轭半径\(OA,OB\)等长。求它们的长度及夹角余弦的绝对值，说明圆化后应同时使用哪些关系。
\par
\solutionanswer
将两共轭半径写成\(A=(3c,2s),B=(-3s,2c)\)，其中\(c^2+s^2=1\)。等长要求\(c^2=s^2=1/2\)，两长度均为\(\sqrt{13/2}\)。数量积为\(-5cs\)，故夹角余弦绝对值为\(5/13\)。圆化后须同时保留共轭半径的垂直关系与原图度量\(9u^2+4v^2\)，才能得到长度及夹角。


% ITEM: quadratic-exercise-27-05
\solutionitem{quadratic-exercise-27-05}{习题4.27（5）}
点\(P\)在椭圆\(x^2/9+y^2/4=1\)上，\(T=(0,1)\)。求\(|PT|\)的最大值、最小值及取等点。可先圆化，将距离的平方化为一个坐标的二次式；比较本题与点到定直线距离问题的区别。
\par
\solutionanswer
作圆化\(x=3u,y=2v\)，有\(u^2+v^2=1\)。将原图距离按两个方向的倍率还原：
\[
 |PT|^2=9u^2+(2v-1)^2=10-5v^2-4v,\qquad -1\le v\le1.
\]
二次式在\(v=-2/5\)处取最大值\(54/5\)，两端值为\(1,9\)，所以
\[
 |PT|_{\max}=\sqrt{54/5},\quad P=\left(\pm\frac{3\sqrt{21}}5,-\frac45\right);
 \qquad |PT|_{\min}=1,\quad P=(0,2).
\]
点到定直线的距离可以化为一个线性式的绝对值；本题的原图距离平方在圆化后仍是加权的二次式。利用圆方程消去\(u^2\)，才把它降为一个坐标的二次函数。


% ITEM: quadratic-exercise-27-06
\solutionitem{quadratic-exercise-27-06}{习题4.27（6）}
椭圆\(x^2/9+y^2/4=1\)上两条共轭半径为\(OA,OB\)。证明\(|OA|^2+|OB|^2=13\)，求\(|OA|\,|OB|\)的取值范围，并指出两端值对应的方向。说明共轭关系与欧氏长度怎样在圆化后配合使用。
\par
\solutionanswer
圆化后的两半径垂直，可取方向 \((c,s),(-s,c)\)，其中 \(c^2+s^2=1\)。还原为 \(A=(3c,2s), B=(-3s,2c)\)，即得 \(OA^2+OB^2=13\)。设两半径夹角为 \(\varphi\)。圆中两半径张成的正方形还原后面积为 \(6\)，于是
\[
  6=OA\cdot OB\sin\varphi\ \le\ OA\cdot OB\ \le\ (OA^2+OB^2)/2=13/2 .
\]
下界要求原图中也垂直，即 \(cs=0\)，两方向为坐标轴；上界要求原图中等长，即 \(c^2=s^2\)，圆化后的方向与坐标轴成 \(45^\circ\)。连续转动两半径，故乘积取遍 \([6,13/2]\)。


% ITEM: quadratic-exercise-27-07
\solutionitem{quadratic-exercise-27-07}{习题4.27（7）}
设\(A=(-3,0)\)、\(B=(3,0)\)，点\(P\)在椭圆\(x^2/9+y^2/4=1\)上。求\(|PA|\,|PB|\)的最大值及取等点。若先将椭圆化为单位圆，能否直接将像图中的距离乘积乘以一个固定常数还原？说明理由。
\par
\solutionanswer
在单位圆中，\(A'B'\) 是直径。对异于两端的 \(P'\)，两弦 \(P'A',P'B'\) 垂直，它们的长度平方和为 \(4\)，故长度乘积至多为 \(2\)。

还原椭圆时，这两条弦各自的伸长倍率，就是沿其方向的一条单位半径的像长。两方向垂直，由上题，两倍率的乘积至多为 \(13/2\)。故 \(|PA|\,|PB|\le 2\cdot13/2=13\)。圆中两弦等长时，\(P'=(0,\pm1)\)，两弦恰沿 \(45^\circ\) 的对角线方向，两个上界同时取到。因此原图的取等点恰为 \((0,\pm2)\)。

这也回答了第二问：换一个非顶点，使两弦不沿对角线方向，两倍率的乘积就严格小于 \(13/2\)，所以不能用一个固定常数还原。


% ITEM: quadratic-exercise-27-08
\solutionitem{quadratic-exercise-27-08}{习题4.27（8）}
设\(P=(x,y)\)在椭圆\(x^2/9+y^2/4=1\)上。求\(|x+y|+|x-y|\)的最大值、最小值及取等点。先研究绝对值式与\(\max\{|x|,|y|\}\)的关系，再比较直接用原图与圆化处理的方便程度。
\par
\solutionanswer
由实数绝对值的分情况关系，原式等于\(2\max(|x|,|y|)\)。最大值为\(6\)，取于\((\pm3,0)\)。设\(r=\max(|x|,|y|)\)，则\(1=x^2/9+y^2/4\le13r^2/36\)，故最小值为\(12/\sqrt{13}\)，在\((\pm6/\sqrt{13},\pm6/\sqrt{13})\)四点取到。若令 \(x=3u,\ y=2v\)，约束虽化为单位圆，目标却成为 \(2\max(3|u|,2|v|)\)，仍要比较两个带不同系数的量。这里直接在原图中以 \(r\) 同时控制两坐标更短。


% ITEM: quadratic-exercise-27-09
\solutionitem{quadratic-exercise-27-09}{习题4.27（9）}
椭圆\(x^2/9+y^2/4=1\)的内接矩形各边平行于坐标轴。求其周长的最大值及取等顶点，再求面积的最大值。取得最大周长与最大面积的矩形是否相同？联系线性式的最值解释。
\par
\solutionanswer
取第一象限顶点\((x,y)\)，周长为\(4(x+y)\)，最大\(4\sqrt{13}\)，顶点为\((\pm9/\sqrt{13},\pm4/\sqrt{13})\)。面积\(4xy\)最大\(12\)，在\((\pm3/\sqrt2,\pm\sqrt2)\)取到。两组图形不同：前者是线性式的最值，后者是坐标乘积的最值。


% ITEM: quadratic-exercise-27-10
\solutionitem{quadratic-exercise-27-10}{习题4.27（10）}
椭圆\(x^2/a^2+y^2/b^2=1\)（\(a>b>0\)）上，\(x+y\)的最大值为\(\sqrt{13}\)，\(2x+y\)的最大值为\(2\sqrt{10}\)。求椭圆方程，并求上述两个线性式各自取得最大值的点。
\par
\solutionanswer
由线性式的椭圆最大值公式，\(a^2+b^2=13,4a^2+b^2=40\)，解得\(a^2=9,b^2=4\)。方程为\(x^2/9+y^2/4=1\)。两个取等点依次为\((9/\sqrt{13},4/\sqrt{13})\)、\((9/\sqrt{10},2/\sqrt{10})\)。


% ITEM: quadratic-exercise-27-11
\solutionitem{quadratic-exercise-27-11}{习题4.27（11）}
设\(A,B\)是椭圆\(x^2/9+y^2/4=1\)上两点，满足\(|OA|=|OB|\)、\(\angle AOB=60^\circ\)。求\(\triangle OAB\)面积的所有可能值。利用原图中的对称性研究等长条件，再与圆化后的度量表达比较。
\par
\solutionanswer
先在原图中利用等长与椭圆方程，得到两点横坐标平方相同、纵坐标平方相同。夹角为\(60^\circ\)排除相同点和相反点，只剩两种镜像关系。若\(A=(x,y),B=(x,-y)\)，夹角条件给出\(x^2=3y^2\)，面积为\(12\sqrt3/7\)；若\(B=(-x,y)\)，则\(y^2=3x^2\)，面积为\(36\sqrt3/31\)。

圆化后取\(A'=(u,v)\)。第一种情形\(B'=(u,-v)\)，原夹角的条件是
\[
 \frac{9u^2-4v^2}{9u^2+4v^2}=\frac12,
\]
即\(9u^2=12v^2\)；第二种得到\(4v^2=27u^2\)。结合\(u^2+v^2=1\)，并把圆中三角形面积乘\(6\)，仍得到上述两值。圆化后不能直接把\(A',B'\)的普通夹角取为\(60^\circ\)；此题先用原图对称性更简洁。


% ITEM: quadratic-example-045
\solutionitem{quadratic-example-045}{例4.45}
（2018年天津卷文科第19题）椭圆\(x^2/a^2+y^2/b^2=1\)（\(a>b>0\)）的右顶点为\(A\)，上顶点为\(B\)，离心率为\(\sqrt5/3\)，且\(|AB|=\sqrt{13}\)。


（1）求椭圆方程。
\par



（2）直线\(l:y=kx\)（\(k<0\)）交椭圆于\(P,Q\)，交\(AB\)于\(M\)，点\(P,M\)均在第四象限。若\(S_{\triangle BPM}=2S_{\triangle BPQ}\)，求\(k\)。
\par
\solutionanswer
（1）由 \(b^2/a^2=4/9,\ a^2+b^2=13\)，得 \(a=3,b=2\)，椭圆方程为 \(x^2/9+y^2/4=1\)。

（2）以 \(PM,PQ\) 为底，两个三角形等高，故 \(PM=2PQ=4OP\)。\(P,M\) 在同一射线上，所以 \(OM=5OP\)。现在令 \(u=x/3,\ v=y/2\)：\(P\) 的像在单位圆上，\(M\) 的像就在半径 \(5\) 的圆上，而直线 \(AB\) 成为 \(u+v=1\)。因此在 \(M\) 处有
\[
  u+v=1,\quad u^2+v^2=25,\quad uv=-12 .
\]
第四象限给出 \(u=4,v=-3\)。还原即 \(M=(12,-6)\)，故 \(k=-1/2\)。


% ITEM: quadratic-exercise-28-01
\solutionitem{quadratic-exercise-28-01}{习题4.28（1）}
（据2012年湖北卷理科第21题）过单位圆上一点\(A\)作\(x\)轴的垂线，垂足为\(D\)。在这条垂线上取点\(M\)，使\(|DM|=m|DA|\)，其中\(m>0\)、\(m\ne1\)。求\(M\)的轨迹\(C\)及焦点。过原点、斜率为\(k>0\)的直线交\(C\)于\(P,Q\)，\(P\)在第一象限，\(N\)为\(P\)在\(y\)轴上的射影，\(QN\)再交\(C\)于\(H\)。是否存在\(m\)，使得任意\(k>0\)都有\(PQ\perp PH\)？比较第二问留在原坐标与圆化后还原垂直条件的处理。
\par
\solutionanswer
由 \(x_M=x_A,\ |y_M|=m|y_A|\)，轨迹为 \(x^2+y^2/m^2=1\)。\(m>1\) 时焦点为 \((0,\pm\sqrt{m^2-1})\)，\(0<m<1\) 时为 \((\pm\sqrt{1-m^2},0)\)。

圆化后 \(P'Q'\) 是直径，故 \(P'H'\perp Q'H'\)。由 \(P=(p,kp),Q=(-p,-kp),N=(0,kp)\)，可知 \(QH\) 的斜率为 \(2k\)。纵坐标除以 \(m\) 后，此斜率变成 \(2k/m\)，故 \(P'H'\) 的斜率为 \(-m/(2k)\)。还原得 \(k_{PH}=-m^2/(2k)\)，所以 \(k_{PQ}k_{PH}=-m^2/2\)。要对一切 \(k>0\) 垂直，充要条件是 \(m=\sqrt2\)。

留在原坐标则可沿用习题4.24（4）的推导：将 \(H,P\) 的椭圆方程相减，分解两组平方差，就得到从 \(H\) 连向同一直径两端，斜率积为 \(-m^2\)，而 \(k_{QH}=2k\)，故仍有 \(k_{PH}=-m^2/(2k)\)。


% ITEM: quadratic-exercise-28-02
\solutionitem{quadratic-exercise-28-02}{习题4.28（2）}
设\(P\)为椭圆\(x^2/9+y^2/4=1\)外一点。两条过\(P\)的直线分别交椭圆于\(A,B\)和\(C,D\)，其单位方向分别为\(d_1,d_2\)。写出\(|PA|\,|PB|\)与\(|PC|\,|PD|\)之比。两条直线方向互相垂直时，这两个乘积之和是否必为定值？
\par
\solutionanswer
方向幂给出两乘积分别为\((q(P)-1)/q(d_1)\)、\((q(P)-1)/q(d_2)\)，故比值为\(q(d_2)/q(d_1)\)。即使两方向垂直，尽管\(q(d_1)+q(d_2)=1/9+1/4\)固定，两者之积一般仍变化，所以两距离乘积之和不必固定。可取\(P=(0,\sqrt5)\)：方向角\(45^\circ,135^\circ\)给出和\(36/13\)；方向角\(30^\circ,120^\circ\)给出和\(624/217\)。这两对直线均为割线。


% ITEM: quadratic-exercise-28-03
\solutionitem{quadratic-exercise-28-03}{习题4.28（3）}
设\(A=(3,0)\)，\(B=(0,2)\)，\(P\)在椭圆\(x^2/9+y^2/4=1\)上。分别求\(S_{\triangle ABP}\)、\(|OP|\)的最大值。判断这两个目标是否适合用同一种方法处理。
\par
\solutionanswer
圆化后面积由弦\(u+v=1\)与单位圆点的距离决定，最大\(S_{ABP}=3(1+\sqrt2)\)，在\(P=(-3/\sqrt2,-2/\sqrt2)\)取到。\(|OP|\)最大为\(3\)，在\((\pm3,0)\)取到。面积适合直接圆化；欧氏距离宜用原图半轴或还原度量计算，不能忽略长度倍率。


% ITEM: quadratic-exercise-28-04
\solutionitem{quadratic-exercise-28-04}{习题4.28（4）}
设\(A,B\)是椭圆\(x^2/a^2+y^2/b^2=1\)上的点，且\(OA\perp OB\)。求\(|OA|^{-2}+|OB|^{-2}\)，并求\(S_{\triangle OAB}\)的取值范围。比较按原图中的方向设坐标与直接化椭为圆两种处理。
\par
\solutionanswer
取两条垂直半径的单位方向为 \((c,s),(-s,c)\)，长度为 \(r_1,r_2\)。分别代入椭圆，两式相加得
\[
 r_1^{-2}+r_2^{-2}=(c^2/a^2+s^2/b^2)+(s^2/a^2+c^2/b^2)=a^{-2}+b^{-2}.
\]
由 \(2/(r_1r_2)\le r_1^{-2}+r_2^{-2}\)，面积 \(r_1r_2/2\) 至少为 \(a^2b^2/(a^2+b^2)\)，沿两条 \(45^\circ\) 方向时两半径等长，达到下界。

上界用圆化更快：单位圆内两半径所成三角形的面积至多 \(1/2\)，还原后至多 \(ab/2\)，沿两坐标轴即取等。转动这对垂直方向，面积取遍两界之间。上界处两像半径确实垂直，但一般位置的像半径未必垂直，不能据此求下界。


% ITEM: quadratic-exercise-28-05
\solutionitem{quadratic-exercise-28-05}{习题4.28（5）}
求椭圆\(x^2/9+y^2/4=1\)内接三角形面积的最大值，并给出一组取等顶点。取得最大面积时，原图中的三角形是否必须为等边三角形？
\par
\solutionanswer
将椭圆化为单位圆。在面积最大的内接三角形中，固定一边，第三顶点必是距该边所在直线最远的圆上点，因而在这条边的中垂线上。对另两边也如此，故最大三角形必为正三角形，其面积为 \(3\sqrt3/4\)。还原时面积乘 \(6\)，最大值为 \(9\sqrt3/2\)。例如取顶点 \((3,0),(-3/2,\sqrt3),(-3/2,-\sqrt3)\) 即达上界；这三个点在原图中并不构成等边三角形。


% ITEM: quadratic-exercise-28-06
\solutionitem{quadratic-exercise-28-06}{习题4.28（6）}
（据2024年新高考II卷第19题）\(P_1=(5,4)\) 在双曲线 \(x^2-y^2=m\) 上。过 \(P_n=(x_n,y_n)\) 作斜率为 \(k\in(0,1)\) 的直线交左支于 \(Q_n\)，令 \(P_{n+1}\) 为 \(Q_n\) 关于 \(y\) 轴的对称点。当 \(k=1/2\) 时求 \(P_2\)；对一般的 \(k\)，证明 \(\{x_n-y_n\}\) 是公比为 \((1+k)/(1-k)\) 的等比数列，且连续三点 \(P_n,P_{n+1},P_{n+2}\) 所成三角形的面积不变。
\par
\solutionanswer
有 \(m=25-16=9\)。取渐近线坐标 \(u=x-y,\ v=x+y\)，双曲线变为 \(uv=9\)。弦 \(P_nQ_n\) 在新坐标中的斜率为 \(h=(1+k)/(1-k)\)。若 \(P_n=(u_n,v_n)\)，将弦式 \(v=v_n+h(u-u_n)\) 代入 \(uv=9\)，两根之积为 \(-9/h\)，故 \(Q_n=(-v_n/h,-hu_n)\)。

关于原 \(y\) 轴的反射在新坐标中为 \((u,v)\mapsto(-v,-u)\)。因此每一步都是同一个伸缩：
\[
    (u_{n+1},v_{n+1})=(hu_n,v_n/h).
\]
当 \(k=1/2\) 时，\((u_1,v_1)=(1,9)\) 变为 \((3,3)\)，即原坐标 \(P_2=(3,0)\)。一般地，\(x_n-y_n=u_n\) 的公比为 \(h\)。这个伸缩的面积倍率为 \(h\cdot h^{-1}=1\)，且把每个三角形变成下一个；换回原坐标时面积又都乘同一倍率，故 \(S_{n+1}=S_n\)。


% ITEM: quadratic-example-046
\solutionitem{quadratic-example-046}{例4.46 *}
对\(2x^2+2xy+y^2=1\)，取\(u=x,v=x+y\)。写出逆变换矩阵\(S\)及度量矩阵\(G\)，分别求像圆上两个坐标方向还原后的长度倍率、夹角与面积倍率。用这一表达重新检查本节第一个中点轨迹例题。
\solutionanswer
逆变换为\((x,y)=(u,v-u)\)，所以\(S=\begin{bmatrix}1&0\\-1&1\end{bmatrix},G=\begin{bmatrix}2&-1\\-1&1\end{bmatrix}\)。两个坐标方向还原后长度倍率为\(\sqrt2,1\)，夹角余弦为\(-1/\sqrt2\)，即\(3\pi/4\)，面积倍率为\(1\)。中点运算无需度量：圆中\((u-1/2)^2+v^2=1/4\)还原为\((x-1/2)^2+(x+y)^2=1/4\)。首例要求\(B\ne A\)，故中点不能是\(A\)：圆中须去掉\((1,0)\)，还原后须去掉\((1,-1)\)。


% ITEM: quadratic-exercise-29-01
\solutionitem{quadratic-exercise-29-01}{习题4.29（1） *}
（2012年福建卷理科第21题）曲线\(2x^2+2xy+y^2=1\)在矩阵\(A=\left[\begin{smallmatrix}a&0\\b&1\end{smallmatrix}\right]\)（\(a>0\)）对应的变换下成为单位圆。求\(a,b\)及\((A^2)^{-1}\)，再写出圆化坐标中的原度量矩阵。
\par
\solutionanswer
要求\(|AX|^2=2x^2+2xy+y^2\)，比较系数得\(b=1,a^2+b^2=2\)，由\(a>0\)得\(a=1\)。所以
\[
A=\begin{bmatrix}1&0\\1&1\end{bmatrix},\quad(A^2)^{-1}=\begin{bmatrix}1&0\\-2&1\end{bmatrix},\quad
G=(A^{-1})^TA^{-1}=\begin{bmatrix}2&-1\\-1&1\end{bmatrix}.
\]


% ITEM: quadratic-exercise-29-02
\solutionitem{quadratic-exercise-29-02}{习题4.29（2） *}
证明若两个可逆矩阵\(S_1,S_2\)满足\(S_1S_1^T=S_2S_2^T\)，则\(S_2^{-1}S_1\)保持长度。解释同一椭圆的不同圆化表示之间的关系。
\par
\solutionanswer
令\(R=S_2^{-1}S_1\)，则\(RR^T=S_2^{-1}S_1S_1^TS_2^{-T}=I\)，故\(R\)正交并保持长度。同一椭圆的两种圆化坐标之间，只相差单位圆上的正交变换，即旋转或镜像。


% ITEM: quadratic-exercise-29-03
\solutionitem{quadratic-exercise-29-03}{习题4.29（3） *}
设\(H\)正定，\(q(X)=X^THX\)。沿非零方向\(V\)的直线\(X=P+tV\)交\(q(X)=1\)于两个点，参数为\(t_1,t_2\)。证明\(t_1t_2q(V)=q(P)-1\)，并与式\textup{(4.93)}比较。
\par

\par\smallskip\noindent 正文所引关系为：
\begin{equation*}
\overline{PA}\,\overline{PB}
    \left(\frac{r^2}{a^2}+\frac{s^2}{b^2}\right)=q(P)-1 .
\end{equation*}
\solutionanswer
把直线代入，得\(q(V)t^2+2b(P,V)t+q(P)-1=0\)。根与系数关系立即给出\(t_1t_2q(V)=q(P)-1\)。若\(V\)为单位方向，参数就是有向长度；一般方向应将参数乘以\(|V|\)，得到\(|PA|\,|PB|=|q(P)-1||V|^2/q(V)\)，与方向幂一致。


% ITEM: quadratic-exercise-29-04
\solutionitem{quadratic-exercise-29-04}{习题4.29（4） *}
设\(S^THS=I\)、\(G=S^TS\)。证明\(H^{-1}=SS^T\)，且\(G\)与\(H^{-1}\)有相同的特征值。由此解释圆化坐标中的最大、最小长度倍率与原椭圆半轴长的关系。
\par
\solutionanswer
由\(S^THS=I\)，反解得\(H=S^{-T}S^{-1}\)，故\(H^{-1}=SS^T\)。又\(SGS^{-1}=SS^T\)，两矩阵相似，特征值相同。它们的特征值为两半轴平方\(a^2,b^2\)，故圆化后还原向量的最大、最小长度倍率为\(a,b\)。


% ITEM: quadratic-exercise-29-05
\solutionitem{quadratic-exercise-29-05}{习题4.29（5） *}
原平面中的单位圆在坐标变换\(X=SU\)后成为怎样的曲线？用\(G=S^TS\)写出方程，求新坐标平面中该曲线所围区域的面积，并与面积倍率核对。
\par
\solutionanswer
原单位圆满足\(X^TX=1\)，代入得\(U^TGU=1\)，是新坐标中的椭圆。面积为\(\pi/\sqrt{\det G}=\pi/|\det S|\)，还原时再乘\(|\det S|\)，恰恢复原圆面积\(\pi\)。


% ITEM: quadratic-exercise-29-06
\solutionitem{quadratic-exercise-29-06}{习题4.29（6） *}
证明：若椭圆的每一对共轭方向都互相垂直，则该椭圆为圆。分别用曲线矩阵\(H\)与圆化后的度量矩阵\(G\)解释这一结论。
\par
\solutionanswer
沿\(H\)的两个正交特征方向建系，记特征值为\(\lambda_1,\lambda_2\)。取\(d=(1,1),e=(\lambda_2,-\lambda_1)\)，则\(d^THe=0\)，题设迫使\(d\cdot e=\lambda_2-\lambda_1=0\)。故\(H\)为单位矩阵的正数倍，曲线为圆。圆化后同一要求是每对普通正交方向在\(G\)度量下仍正交。先取两个坐标方向，得\(g_{12}=0\)；再取\((1,1),(1,-1)\)，得\(g_{11}-g_{22}=0\)。于是\(G=g_{11}I\)，还原时所有方向的长度倍率相同，单位圆仍为圆。
\par
